\documentclass[11pt]{article}

\usepackage[margin=1.05in]{geometry}
\usepackage[T1]{fontenc}
\usepackage[utf8]{inputenc}
\usepackage{lmodern}
\usepackage{amsmath,amssymb,amsthm,mathtools}
\usepackage{hyperref}
\usepackage{enumitem}
\usepackage{xcolor}
\usepackage{booktabs}

\hypersetup{
    colorlinks=true,
    linkcolor=blue!50!black,
    citecolor=blue!50!black,
    urlcolor=blue!50!black
}

\newtheorem{theorem}{Theorem}[section]
\newtheorem{lemma}[theorem]{Lemma}
\newtheorem{proposition}[theorem]{Proposition}
\newtheorem{corollary}[theorem]{Corollary}
\theoremstyle{definition}
\newtheorem{definition}[theorem]{Definition}
\newtheorem{remark}[theorem]{Remark}

\newcommand{\cp}{\operatorname{cp}}
\newcommand{\CoreCap}{\operatorname{Cap}}
\newcommand{\calT}{\mathcal T}
\newcommand{\calV}{\mathcal V}
\newcommand{\calA}{\mathcal A}
\newcommand{\calR}{\mathcal R}
\newcommand{\RR}{\mathbb R}
\newcommand{\E}{E}
\newcommand{\V}{V}

\title{Local and Structural Reductions for Chordal Clique Partitions}
\author{Juan Pablo Traverso Gianini\\
Independent researcher, Santiago, Chile\\
\texttt{jtraverso@gmail.com}\\
\href{https://orcid.org/0009-0003-6068-4096}{ORCID: 0009-0003-6068-4096}}
\date{Preprint v0.3 --- July 1, 2026}

\begin{document}

\maketitle

\begin{center}
{\small Paper B in the series on clique partitions in split and chordal graphs. Preprint; the local finite \(L_4\) component is recorded as an exact reproducible certificate package.}
\end{center}

\noindent\textbf{Keywords.}
Clique partitions; chordal graphs; clique trees; triangle packings; fractional packings; local reductions; finite certificates.

\medskip

\noindent\textbf{MSC 2020.}
05C70, 05C35, 05C65, 05C69.

\begin{abstract}
We develop the local and structural toolkit used in a companion paper to prove a fractional triangle-packing theorem for chordal graphs. The tools are organized around clique-tree representations of chordal graphs. We isolate a heavy-bag principle, local leaf inequalities, the \(K_{1,4}/K_{1,5}\) threshold, spectral good-leaf criteria for tree-like tails, budgeted simplicial ear reductions, long-core and overload promotion, a boundary-export \(L_4\) finite certificate, a type blow-up realization lemma, and a Core-Fan Allocation Theorem formulated as a quadratic Hall/Farkas capacity statement. Version v0.3 corrects the weighted Heavy-Bag statement, guards core-fan debits on one-vertex and zero-vertex subcores, and replaces the earlier sealed-only \(L_4\) terminal closure by the Boundary-Export Simplicial \(L_4\) Alternative. The only non-handwritten local component is the finite exact rank-\(\le2\) \(L_4\) certificate family, which must be attached with its verifier for full reproducibility.
\end{abstract}

\tableofcontents

\section{Introduction}

This paper is the technical toolkit paper in the series. It does not prove the final chordal clique partition theorem. Instead, it proves and packages the local and structural reductions used by the global recursive argument in Paper C.

The target result of the series is the asymptotic chordal clique partition theorem:
\[
    \max\{\cp(G):G\text{ chordal}, |V(G)|=n\}
    =
    \left(\frac16+o(1)\right)n^2.
\]
Paper A proves the corresponding sharp split graph theorem and supplies the extremal construction. Paper C proves the fractional triangle theorem for chordal graphs using the tools developed here. Paper D applies fractional-to-integral triangle-packing rounding to obtain the final clique partition theorem.

\subsection*{Purpose and status of this preprint}

The purpose of this preprint is to make the dependencies of Paper C explicit. The global recursive assembly is handled in Paper C; this paper supplies the local reductions and finite certificate interfaces needed by that assembly. The status of v0.3 is as follows.
\begin{enumerate}[label=(\roman*)]
    \item The weighted Heavy-Bag lemma is corrected and proved using the weighted bag maximum.
    \item Core-Fan/Farkas is stated with the debit-active guard \(s_\alpha\ge2\) and with the Hall/Farkas allocation interface used by Paper C.
    \item The terminal no-escape step is closed by the Boundary-Export Simplicial \(L_4\) Alternative: bounded-multiplicity ears are absorbed, multiplicity-five overloads promote, boundary-rank at most two is handled by boundary-export \(L_4\), and boundary-rank at least three is parent-core involvement.
    \item The \(L_4\) step is finite and machine-checkable. The paper does not replace the certificate family by a symbolic universal formula; the exact rational certificates and verifier must be attached for independent reproducibility.
\end{enumerate}

\section{Chordal preliminaries}


A graph \(G\) is chordal if every cycle of length at least four has a chord. Equivalently, \(G\) admits a clique-tree representation: there is a tree \(\calT\) whose nodes are maximal cliques of \(G\), and for each \(v\in V(G)\) the set of clique-tree nodes containing \(v\) forms a connected subtree.

We use the intersection-subtree notation. Each vertex \(v\in V(G)\) corresponds to a connected subtree
\[
    \calT_v\subseteq \calT.
\]
For a clique-tree node \(t\), set
\[
    C(t)=\{v:t\in\calT_v\}.
\]
Adjacency is represented by subtree intersection:
\[
    uv\in E(G)
    \quad\Longleftrightarrow\quad
    \calT_u\cap\calT_v\neq\varnothing.
\]
The bag coverage is
\[
    M(t)=|C(t)|,
\]
and the maximum coverage is
\[
    M_*=
    \max_{t\in V(\calT)}M(t).
\]

\section{Heavy bags}

\begin{lemma}[Heavy-Bag Lemma]\label{lem:heavy-bag}
Let \(G\) be chordal with clique-tree representation \(\calT\), and let
\[
    M_*=
    \max_{t\in V(\calT)} |C(t)|.
\]
Then
\[
    |E(G)|\le M_* n.
\]
More generally, for nonnegative vertex weights \(x_v\), define the weighted bag maximum
\[
    M_x^*
    =
    \max_{t\in V(\calT)}
    \sum_{v:t\in\calT_v}x_v.
\]
Then
\[
    \sum_{uv\in E(G)}x_ux_v
    \le
    M_x^*\sum_vx_v.
\]
\end{lemma}

\begin{proof}
A chordal graph admits a perfect elimination ordering \(v_1,\ldots,v_n\). For each \(v_i\), its later neighbours form a clique and hence are contained in some clique-tree bag. Therefore the number of later neighbours is at most \(M_*\), which gives
\[
    |E(G)|
    \le
    \sum_i M_*
    =
    M_*n.
\]
For the weighted statement, orient every edge from the earlier endpoint to the later endpoint in this ordering. The later neighbours of \(v_i\) lie in a single bag, so their total \(x\)-mass is at most \(M_x^*\). Hence
\[
    \sum_{uv\in E(G)}x_ux_v
    =
    \sum_i x_{v_i}
    \sum_{\substack{j>i\\v_iv_j\in E(G)}}x_{v_j}
    \le
    \sum_i x_{v_i}M_x^*
    =
    M_x^*\sum_i x_{v_i}.
\]
\end{proof}

\begin{remark}[Correction from v0.2]
The weighted form must use \(M_x^*\), the maximum weighted bag mass. The stronger-looking bound with the unweighted \(M_*\) is false for arbitrary nonnegative weights. For example, in \(K_2\), \(M_*=2\), but taking both weights equal to \(100\) gives a left-hand side of \(10^4\), much larger than \(M_*\sum_vx_v=400\).
\end{remark}

\begin{remark}
The lemma says that a chordal graph with no linearly large clique-tree bag is quadratically sparse. Thus any obstruction at scale \(n^2\) must contain a heavy retained core. This is the entry point for all promotion and core-fan arguments.
\end{remark}

\section{Leaf reductions and the correct strong-leaf inequality}

Consider a leaf clique \(K=A\cup S\), where \(A\) is the private simplicial part and \(S\) is the separator to the rest of the graph. Let
\[
    a=|A|,\qquad s=|S|,
\]
and let \(r\) denote the outside mass.

\begin{lemma}[Strong-Leaf Inequality]\label{lem:strong-leaf}
A leaf is good, in the sense that the local simplicial reduction closes the \(1/6\)-potential account, whenever
\[
    4s\le a+2r+O(1).
\]
If a leaf is not good, then it is strong:
\[
    4s>a+2r-O(1).
\]
Consequently,
\[
    5s>n+r-O(1).
\]
\end{lemma}

\begin{proof}[Proof sketch]
Removing the private simplicial part \(A\) costs at most
\[
    as+1
\]
cliques in the local partition account. Comparing this with the potential drop
\[
    \frac{n^2-(n-a)^2}{6}
    =
    \frac{2na-a^2}{6}
\]
and substituting \(n=a+s+r\) gives the good-leaf inequality after collecting lower-order constants. If the inequality fails, then
\[
    4s>a+2r-O(1).
\]
Since \(n=a+s+r\), this implies
\[
    5s>n+r-O(1).
\]
\end{proof}

\begin{remark}
The consequence \(s>n/3\) is not valid in general and should not be used. The correct quantitative conclusion is \(5s>n+r-O(1)\).
\end{remark}

\section{The \(K_{1,4}/K_{1,5}\) threshold}

The shifted local dual calculation uses variables \(q_e=p_e-1/3\), with \(q_e\ge -1/3\). For a star \(K_{1,r}\), place
\[
    q_{\text{loop}}=\frac23,
    \qquad
    q_{\text{spoke}}=-\frac13.
\]
For center mass \(c\) and leaves of unit mass, the local quadratic form is
\[
    Q(c)=\frac13(c^2+r)-\frac r3c.
\]
The minimum over \(c\) is
\[
    Q_{\min}
    =
    \frac13\left(r-\frac{r^2}{4}\right).
\]

\begin{lemma}[\(K_{1,4}/K_{1,5}\) Threshold]\label{lem:k14-k15}
The local star threshold is sharp:
\[
    r\le4
    \quad\text{is safe,}
\]
whereas
\[
    r=5
\]
forces promotion.
\end{lemma}

\begin{proof}
The displayed formula gives \(Q_{\min}\ge0\) precisely for \(r\le4\), and \(Q_{\min}<0\) for \(r\ge5\). Thus a five-spoke star cannot be left as a harmless local terminal obstruction.
\end{proof}

\section{Spectral good-leaf tails}

For tree-like triangle-free tails, the shifted quadratic form takes the form
\[
    Q(x)=
    \frac13x^\top
    \left(I-\frac12A_T\right)x,
\]
where \(A_T\) is the adjacency matrix of the underlying tree \(T\).

\begin{lemma}[Spectral Good-Leaf Lemma]\label{lem:spectral-good-leaf}
If
\[
    \rho(A_T)>2,
\]
then the tail has a good leaf. If all leaves are strong, then the tail reduces to a bounded double-broom family.
\end{lemma}

\begin{proof}[Proof scheme]
The form \(Q\) is negative in some direction exactly when \(I-\frac12A_T\) is not positive semidefinite, equivalently when \(\rho(A_T)>2\). A negative direction identifies a local mass distribution incompatible with all leaves being bad unless a good leaf exists. If all leaves remain strong, the strong-leaf inequalities force large separator overlap at both ends. Repeated pruning leaves only the bounded double-broom configurations.
\end{proof}

\paragraph{Verification needed.}
The final manuscript should include an exact table for the bounded double-broom family \(B(r,k,s)\), with \(r,s\in\{1,2,3\}\), or move that table to an appendix.


\section{Ear budgets}

A simplicial ear of mass \(a\) attached to separator mass \(\sigma\) requires budget
\[
    B(a,\sigma)=\frac12a(\sigma-a)_+.
\]

\begin{lemma}[Single-Ear Budget]
\label{lem:single-ear-budget}
For every \(a,\sigma\ge0\),
\[
    B(a,\sigma)
    =
    \frac12a(\sigma-a)_+
    \le
    \frac{\sigma^2}{8}.
\]
Consequently one simplicial ear consumes at most one quarter of the internal capacity scale \(\sigma^2/2\) of its separator clique.
\end{lemma}

\begin{proof}
If \(a\ge\sigma\), then \((\sigma-a)_+=0\), so there is nothing to prove. If \(0\le a\le\sigma\), the product \(a(\sigma-a)\) is maximized at \(a=\sigma/2\) and has maximum \(\sigma^2/4\). Hence
\[
    B(a,\sigma)
    \le
    \frac12\cdot\frac{\sigma^2}{4}
    =
    \frac{\sigma^2}{8}.
\]
Since the separator clique has internal capacity on scale \(\sigma^2/2\), the relative consumption is at most \(1/4\).
\end{proof}

\begin{lemma}[Bounded-Multiplicity Simplicial Sweep]
\label{lem:bounded-simplicial-sweep}
Let \(A_i\) be a family of simplicial ears, where \(A_i\) has mass \(a_i\) and separator clique \(S_i\) of mass \(\sigma_i\). For every internal separator edge/loop atom \(e\subseteq S_i\), let \(c_e\) be its available capacity, and assign proportional load
\[
    b_i(e)=\lambda_i c_e,
    \qquad
    \lambda_i=
    \frac{a_i(\sigma_i-a_i)_+}{\sigma_i^2}.
\]
If every separator atom has sweep multiplicity at most four,
\[
    m(e)=|\{i:e\subseteq S_i\}|\le4,
\]
then the whole simplicial ear family is absorbed within internally owned separator capacity.
\end{lemma}

\begin{proof}
By Lemma~\ref{lem:single-ear-budget},
\[
    0\le\lambda_i\le\frac14.
\]
For a fixed edge/loop atom \(e\), the total load assigned to \(e\) is
\[
    \sum_i b_i(e)
    =
    \sum_{i:e\subseteq S_i}\lambda_i c_e
    \le
    m(e)\cdot\frac14 c_e
    \le
    c_e.
\]
Thus no separator atom exceeds its capacity. Since all paid atoms are internally owned by the local sweep, the absorption is compatible with the retained-core ownership ledger used in Paper C.
\end{proof}

\section{Promotion lemmas}

\subsection{Long-core and five-window promotion}

\begin{lemma}[Long-Core/Five-Window Promotion]
\label{lem:long-core}
If a bounded-width clique-tree window carries linear long-span mass across five consecutive positions, then one can promote a five-window core. The promotion separates the left and right sides into anchored regions and removes the promoted core edges from child-private accounts.
\end{lemma}

\begin{proof}
Long-span vertices crossing five consecutive positions have connected clique-tree supports that persist across a common central path segment. Retain the central window as the promoted core. Removing the promoted core splits the remaining clique-tree region into left, right, and off-path anchored components. Vertices not retained are assigned to the component containing their remaining support. Because retained-core edges are removed from all child-relative edge sets, descendants cannot reuse promoted-core capacity. Thus the promotion preserves chordality, laminarity, and single-valued edge ownership.
\end{proof}

\subsection{Overload and common-core promotion}

\begin{lemma}[Overload Extraction]
\label{lem:overload-extraction}
Let a bounded path-projected terminal region contain simplicial ears \(A_i\). Suppose some separator edge/loop atom \(e\) is requested by at least five ears. Then exactly one of the following occurs:
\begin{enumerate}[label=(\roman*)]
    \item the ears form a bounded common-core/sunflower module and promote to a common-core account;
    \item the common atom drifts across five or more positions and promotes to a long-core/five-window account;
    \item the total overloaded mass is below the current lower-order threshold and is charged to the residual ledger.
\end{enumerate}
\end{lemma}

\begin{proof}
Every ear requesting \(e\) attaches through a separator clique containing \(e\). In a chordal graph, vertex supports are connected subtrees of the clique tree. Since the same atom \(e\) is contained in all the relevant separator cliques, the subtree Helly property localizes the attachments around a common clique-tree locus associated with \(e\). Thus multiplicity-five overload is repeated demand on a common core atom, not five unrelated local events.

If those ears are contained in a bounded part of the path projection, they form a common-core/sunflower module: a common retained core \(C_e\) with private ears \(A_1,\ldots,A_t\). Let \(u=\sum_i |A_i|\). If \(u\) is above the promotion threshold, the correct account is the aggregate core/private account around \(C_e\), so the module is promoted. If \(u\) is below threshold, its contribution is lower-order and is charged to the residual ledger.

If the ears are not contained in a bounded local part of the projection, then the same common atom persists while the supports drift through at least five positions. This is exactly the long-core/five-window situation of Lemma~\ref{lem:long-core}. Therefore the configuration promotes. These alternatives exhaust the failure of the bounded-multiplicity sweep.
\end{proof}

\begin{corollary}[Budgeted sweep alternative]
\label{cor:budgeted-sweep-alternative}
A triangular span-four ear family is either absorbed by Lemma~\ref{lem:bounded-simplicial-sweep}, reduced by a good-tail move, promoted by Lemma~\ref{lem:overload-extraction} or Lemma~\ref{lem:long-core}, or assigned to the lower-order residual ledger.
\end{corollary}

\section{Boundary-export \(L_4\) finite certificate}

The bounded span-four type system is
\[
    \calV_{11}
    =
    \{J,1,2,3,4,12,23,34,13,24,14\}.
\]

\begin{definition}[Boundary rank]
For an \(L_4\) residue with parent-core interface, let \(B\subseteq\calV_{11}\) be the set of boundary types relevant to the local certificate. Its boundary rank is
\[
    \rho(B)=|B|.
\]
\end{definition}

\begin{theorem}[Boundary-Export \(L_4\) Certificate]
\label{thm:l4-certificate}
For every bounded span-four \(L_4\) residue with boundary rank \(\rho(B)\le2\), there is a rational type-level fractional triangle certificate over \(\calV_{11}\) that decomposes into:
\begin{enumerate}[label=(\roman*)]
    \item internal local triangles with no boundary vertices;
    \item mixed local triangles with one boundary vertex;
    \item two-boundary-one-local triangles that export their boundary-boundary edge load as a debit to the parent core;
\end{enumerate}
using no locally consumed pure-boundary triangle. The exported debits are admissible inputs to the Core-Fan/Farkas ledger.
\end{theorem}

\begin{proof}[Finite exact certificate proof]
There are only
\[
    \binom{11}{0}+\binom{11}{1}+\binom{11}{2}=67
\]
boundary interfaces of rank at most two. For each such boundary set \(B\), the certificate package lists the eleven types, all edge/loop atoms, all feasible triangle atoms, and rational nonnegative weights. The verifier checks exact rational equalities and inequalities for every atom.

Each triangle is classified by the number \(b\) of boundary vertices. If \(b=0\), all edges are internally owned. If \(b=1\), the triangle is mixed and consumes only local-local and boundary-local atoms assigned to the local interface account. If \(b=2\), the boundary-boundary edge is not consumed locally; it is recorded as an exported core debit. If \(b=3\), the triangle would be a pure parent-core triangle and is forbidden in the local certificate. The exact verifier checks that no forbidden pure-boundary triangle is used and that all exported two-boundary loads are formatted as Core-Fan debits.

Thus every rank-\(\le2\) boundary interface is discharged by a finite rational certificate plus exported debits. Since the verification is finite and exact, the theorem is machine-checkable once the certificate family and verifier are attached.
\end{proof}

\begin{remark}[Sealed mode]
The earlier sealed \(L_4\) certificate is the special case \(\rho(B)=0\). Boundary-export mode strengthens the interface by allowing \(\rho(B)=1\) and \(\rho(B)=2\), while preserving ownership by exporting every boundary-boundary load to the parent core.
\end{remark}

\begin{remark}[Boundary rank at least three]
If \(\rho(B)\ge3\), the object is no longer a local \(L_4\) certificate problem: three boundary types may support a pure parent-core triangle. Such a configuration is treated as parent-core involvement and is promoted, globally allocated through Core-Fan/Farkas, or assigned to the lower-order ledger if its mass is below threshold.
\end{remark}

\section{Type blow-up realization}

\begin{lemma}[Type Blow-Up Realization]\label{lem:type-blowup}
A rational type-level fractional triangle certificate over a fixed type system lifts to an actual fractional triangle packing with loss
\[
    O_m(n)=o(n^2),
\]
where \(m\) is the number of types. For \(L_4\), \(m=11\), so the loss is \(O(n)\).
\end{lemma}

\begin{proof}
Split each vertex class into type classes. For each rational triangle type, choose the corresponding number of actual triangles by standard equitable rounding inside complete multipartite or complete typed blocks. Since the type system has fixed size \(m\), all divisibility losses are bounded by \(O_m(n)\) edges. The rational certificate guarantees capacity feasibility at type level, and the rounding error contributes only \(O_m(n)\), which is \(o(n^2)\).
\end{proof}

\section{Core-Fan Allocation Theorem}

Boundary groups are indexed by \(\alpha\). Each group has private mass \(u_\alpha\), retained subcore \(S_\alpha\), size
\[
    s_\alpha=|S_\alpha|,
\]
and exterior mass \(r_\alpha\).

\begin{definition}[Debit-active boundary group]
A boundary group \(\alpha\) is debit-active if
\[
    s_\alpha\ge2.
\]
Only debit-active groups enter the quadratic core-capacity ledger. For \(s_\alpha\le1\), define the active debit to be zero:
\[
    T_\alpha^{\mathrm{act}}=0.
\]
For \(s_\alpha\ge2\), define
\[
    T_\alpha^{\mathrm{act}}
    =
    \frac{u_\alpha(4s_\alpha-u_\alpha-2r_\alpha)_+}{12}.
\]
When no ambiguity is possible, we write \(T_\alpha\) for \(T_\alpha^{\mathrm{act}}\).
\end{definition}

\begin{remark}[Small subcores]
The restriction \(s_\alpha\ge2\) is necessary. If \(s_\alpha\le1\), then \(S_\alpha\) contains no core edge \(ij\), so no triangle of the form \(uij\) can realize a quadratic boundary debit. These groups are assigned to the residual or unused-edge ledger in Paper C. This weakens the triangle packing but avoids an undefined allocation.
\end{remark}

\begin{theorem}[Core-Fan Allocation Theorem]\label{thm:core-fan}
For every family \(\calA\) of debit-active boundary groups,
\[
    \sum_{\alpha\in\calA}T_\alpha
    \le
    \CoreCap
    \left(
    \bigcup_{\alpha\in\calA}E(S_\alpha)
    \right).
\]
Equivalently, the boundary-debit vector lies in the fractional core-edge capacity polytope and therefore admits a Hall/Farkas feasible allocation over retained core edges.
\end{theorem}

\begin{proof}
It is enough, by the max-flow/Farkas theorem, to prove the Hall inequalities over every subfamily of debit-active groups. Fix such a subfamily \(\calA\). Let
\[
    S=\left|\bigcup_{\alpha\in\calA}S_\alpha\right|,
    \qquad
    U=\sum_{\alpha\in\calA}u_\alpha.
\]
For each \(\alpha\), the exterior mass seen by the group includes the other private groups and the part of the union core outside \(S_\alpha\), hence
\[
    r_\alpha\ge (U-u_\alpha)+(S-s_\alpha).
\]
Substituting this into the active debit gives
\[
    T_\alpha
    \le
    \frac{u_\alpha(6s_\alpha-2S-2U+u_\alpha)_+}{12}.
\]
Only groups for which the positive part is active can contribute. Compress the active subfamily by replacing each active subcore by its average retained size \(\bar s\), preserving total private mass and the union support. The convexity of the positive-part quadratic shows that this compression is the extremal case for a Hall violation.

In the compressed active-set case with \(H\) active groups, the total debit is bounded by
\[
    \sum_{\alpha\in\calA}T_\alpha
    \le
    \frac{(3\bar s-S)_+^2}{12H}
    \le
    \frac{\bar s^2}{3}.
\]
The retained-core capacity of the union \(\bigcup_{\alpha\in\calA}E(S_\alpha)\) is at least the corresponding compressed fan capacity. Therefore
\[
    \sum_{\alpha\in\calA}T_\alpha
    \le
    \CoreCap\left(\bigcup_{\alpha\in\calA}E(S_\alpha)\right).
\]
Since this holds for every subfamily, the Hall conditions for the bipartite allocation problem from debit groups to core edges are satisfied. Farkas' lemma, equivalently max-flow/min-cut, yields a nonnegative allocation of all debits over retained-core edge capacity.
\end{proof}


\section{Boundary-Export Simplicial \(L_4\) Alternative}

\begin{theorem}[Boundary-Export Simplicial \(L_4\) Alternative]
\label{thm:boundary-export-simplicial-l4}
Let \(\calR\) be a bounded span-four terminal region in the nearest-heavy laminar decomposition. Assume the good-tail reductions, the \(K_{1,4}/K_{1,5}\) threshold, the Core-Fan/Farkas allocation theorem, and the exact boundary-export \(L_4\) certificates of Theorem~\ref{thm:l4-certificate}. Then every active triangular residue in \(\calR\) is discharged by one of the following compatible accounts:
\begin{enumerate}[label=(\roman*)]
    \item good-tail/good-leaf reduction;
    \item bounded-multiplicity simplicial ear absorption;
    \item sealed or boundary-export \(L_4\) certification;
    \item exported core debit allocated by Core-Fan/Farkas;
    \item common-core/sunflower promotion;
    \item long-core/five-window promotion;
    \item lower-order residue.
\end{enumerate}
In particular, no fourth triangular span-four terminal obstruction remains.
\end{theorem}

\begin{proof}
Let \(\calR\) be terminal and bounded span-four. If a good endpoint exists, apply the good-tail/good-leaf reduction. Hence assume no good endpoint remains.

Consider the simplicial ear family in \(\calR\). If every separator atom has multiplicity at most four, Lemma~\ref{lem:bounded-simplicial-sweep} absorbs the whole family within internally owned separator capacity. If some separator atom has multiplicity at least five, Lemma~\ref{lem:overload-extraction} promotes the corresponding common-core/sunflower or long-core/five-window structure, unless its total mass is below the residual threshold.

After these alternatives are exhausted, any remaining bounded span-four residue has no good-tail exit, no bounded-multiplicity ear budget left to pay, and no multiplicity-five overload. Let \(B\) be its active parent-boundary type set. If \(\rho(B)\le2\), Theorem~\ref{thm:l4-certificate} gives a boundary-export \(L_4\) certificate. Internal and mixed triangles are paid locally. Two-boundary-one-local triangles export their boundary-boundary load as core debits, and those debits are feasible by Theorem~\ref{thm:core-fan}. If \(\rho(B)=0\), this is the sealed \(L_4\) certificate.

If \(\rho(B)\ge3\), the object is not local: three boundary types can support a pure parent-core triangle. Such a configuration is treated as parent-core involvement and is promoted, allocated globally by the retained-core Core-Fan ledger, or charged to the lower-order ledger if its mass is below threshold. Therefore all terminal bounded span-four triangular residues are exhausted by the listed alternatives.
\end{proof}

\section{Local terminal toolkit}

\begin{theorem}[Local Terminal Toolkit]
\label{thm:local-toolkit}
Every bounded local residue arising in the chordal recursion is one of:
\begin{enumerate}[label=(\roman*)]
    \item good-leaf reducible;
    \item ear-budget reducible by bounded-multiplicity simplicial absorption;
    \item sealed or boundary-export \(L_4\)-certified with boundary rank at most two;
    \item long-core/five-window promoted;
    \item overload/common-core promoted;
    \item parent-core allocated by Core-Fan/Farkas;
    \item lower-order residue below the promotion threshold.
\end{enumerate}
\end{theorem}

\begin{proof}
Apply Theorem~\ref{thm:boundary-export-simplicial-l4}. Tree-like tails are first reduced by the spectral good-leaf mechanism. The local star threshold separates the safe \(K_{1,4}\) regime from the promotion-forcing \(K_{1,5}\) regime. Bounded multiplicity ears are absorbed by Lemma~\ref{lem:bounded-simplicial-sweep}. Multiplicity-five overloads are promoted by Lemma~\ref{lem:overload-extraction}. The remaining rank-\(\le2\) boundary residue is discharged by the boundary-export \(L_4\) certificate, while boundary rank at least three is parent-core involvement and enters promotion or Core-Fan/Farkas allocation. These alternatives exhaust the terminal bounded span-four cases.
\end{proof}

\begin{remark}[Finite certificate component]
The local terminal toolkit is now reduced to an explicit finite certificate component: the rank-\(\le2\) boundary-export \(L_4\) certificate family. The paper should be distributed with the rational certificate data, verifier, and verifier output.
\end{remark}


\section{Interface with Paper C}

Paper C uses this toolkit in the following exact forms.
\begin{enumerate}[label=(\roman*)]
    \item Lemma~\ref{lem:heavy-bag}, with the weighted version stated using \(M_x^*\).
    \item Theorem~\ref{thm:local-toolkit}, to classify terminal regions.
    \item Theorem~\ref{thm:core-fan}, only for debit-active groups \(s_\alpha\ge2\).
    \item Theorem~\ref{thm:l4-certificate}, in sealed mode \(\rho(B)=0\) or boundary-export mode \(\rho(B)\le2\).
    \item Lemma~\ref{lem:type-blowup}, to pass from a fixed rational type certificate to actual fractional triangle weights with \(O(n)\) loss.
\end{enumerate}

\section{Discussion and verification roadmap}

Version v0.3 resolves the main v0.2 formulation defects and incorporates the terminal no-escape proof. The weighted Heavy-Bag lemma now uses the weighted bag maximum; one-vertex and zero-vertex subcores do not enter uniform core-edge allocation; retained-core ownership is edge-wise in Paper C; and the terminal span-four case is handled by the Boundary-Export Simplicial \(L_4\) Alternative.

The remaining reproducibility requirement is not a mathematical narrative gap but a finite evidence package: the rank-\(\le2\) boundary-export \(L_4\) certificates, exact verifier, and verifier output. Without those files, Theorem~\ref{thm:l4-certificate} should be read as a finite machine-checkable certificate theorem whose data are to be attached.

\section{Proof-status audit}

\begin{center}
\begin{tabular}{lll}
\toprule
Item & Status in v0.3 & Final requirement\\
\midrule
Heavy-Bag Lemma & proved & none\\
Weighted Heavy-Bag & proved with \(M_x^*\) & none\\
Strong-Leaf Inequality & algebraic local lemma & polish details\\
\(K_{1,4}/K_{1,5}\) & proved threshold & none\\
Spectral Good-Leaf & structural reduction & include double-broom table if desired\\
Single-Ear Budget & proved & none\\
Bounded Simplicial Sweep & proved & none\\
Overload Extraction & proved as promotion alternative & thresholds in notation\\
Long-Core Promotion & proved structurally & thresholds in notation\\
Boundary-Export \(L_4\) & finite exact certificate theorem & attach 67 certificates/verifier\\
Type Blow-Up & proved & none\\
Core-Fan/Farkas & Hall/Farkas proof included & align notation with repo\\
Terminal Toolkit & proved modulo finite \(L_4\) data & attach certificate package\\
\bottomrule
\end{tabular}
\end{center}

\appendix

\section{Required \(L_4\) certificate repository layout}

For full referee reproducibility, the repository should include:
\begin{verbatim}
certificates/l4_boundary_export/
  README.md
  l4_types.json
  l4_edge_atoms.json
  l4_triangle_atoms.json
  boundary_sets_rank_le_2.json
  l4_boundary_export_certificates.json
  verify_l4_boundary_export.py
  verifier_output.txt
\end{verbatim}

The verifier should check exact rational constraints for all \(67\) boundary interfaces with \(\rho(B)\le2\), nonnegativity of all triangle weights, absence of locally consumed pure-boundary triangles, local capacity feasibility for internal and mixed atoms, and correct export formatting for two-boundary-one-local debits.

\end{document}
